\begin{ex}
Để vẽ một hình tròn bán kính 4 cm bằng compa, ta mở hai đầu compa cách nhau bao nhiêu xăng-ti-mét?
\shortans{4}
\loigiai{
Mở compa đúng bằng độ dài bán kính: 4 cm.
}
\end{ex}

\begin{ex}
Để vẽ hình tròn có đường kính 10 cm, ta mở hai đầu compa cách nhau bao nhiêu xăng-ti-mét?
\shortans{5}
\loigiai{
Bán kính = $10 : 2 = 5$ cm.
}
\end{ex}

\begin{ex}
Bạn Mai vẽ 4 hình tròn bằng nhau, mỗi hình có bán kính 3 cm. Tổng đường kính của 4 hình tròn đó là bao nhiêu xăng-ti-mét?
\shortans{24}
\loigiai{
Mỗi hình tròn có đường kính $3 \times 2 = 6$ cm. Tổng đường kính = $6 \times 4 = 24$ cm.
}
\end{ex}

\begin{ex}
Vẽ một hình vuông có cạnh 5 cm trên giấy kẻ ô vuông. Chu vi hình vuông đó là bao nhiêu xăng-ti-mét?
\shortans{20}
\loigiai{
$5 \times 4 = 20$ cm.
}
\end{ex}

\begin{ex}
Vẽ một hình chữ nhật có chiều dài 6 cm, chiều rộng 4 cm. Chu vi hình chữ nhật vẽ được là bao nhiêu xăng-ti-mét?
\shortans{20}
\loigiai{
$(6 + 4) \times 2 = 20$ cm.
}
\end{ex}

\begin{ex}
Dùng 3 que tính bằng nhau có thể xếp thành hình gì? (Hình tam giác ghi 1, tứ giác ghi 2).
\shortans{1}
\loigiai{
3 que tính xếp thành hình tam giác.
}
\end{ex}

\begin{ex}
Dùng 4 que tính bằng nhau có thể xếp thành hình gì có 4 góc vuông? (Hình vuông ghi 1, hình tròn ghi 2).
\shortans{1}
\loigiai{
Xếp thành hình vuông.
}
\end{ex}

\begin{ex}
Cần ít nhất bao nhiêu que tính bằng nhau để xếp thành một hình chữ nhật (không phải hình vuông)?
\shortans{6}
\loigiai{
Mỗi chiều dài 2 que, mỗi chiều rộng 1 que: $2 \times 2 + 1 \times 2 = 6$ que tính.
}
\end{ex}

\begin{ex}
Cần bao nhiêu que tính bằng nhau để ghép thành 2 hình vuông riêng biệt?
\shortans{8}
\loigiai{
$4 \times 2 = 8$ que.
}
\end{ex}

\begin{ex}
Nếu ghép chung 1 cạnh, dùng bao nhiêu que tính bằng nhau để tạo thành 2 hình vuông liền nhau?
\shortans{7}
\loigiai{
$4 + 3 = 7$ que.
}
\end{ex}

\begin{ex}
Bạn Lan dùng dây kim tuyến viền quanh một khung ảnh hình vuông cạnh 15 cm. Bạn Lan cần bao nhiêu xăng-ti-mét dây?
\shortans{60}
\loigiai{
$15 \times 4 = 60$ cm.
}
\end{ex}

\begin{ex}
Một dải ruy băng trang trí hình chữ nhật dài 20 cm, rộng 10 cm. Chu vi dải ruy băng đó là bao nhiêu xăng-ti-mét?
\shortans{60}
\loigiai{
$(20 + 10) \times 2 = 60$ cm.
}
\end{ex}

\begin{ex}
Bạn Nam dùng compa vẽ 3 hình tròn đồng tâm có bán kính lần lượt là 2 cm, 4 cm, 6 cm. Đường kính của hình tròn lớn nhất là bao nhiêu xăng-ti-mét?
\shortans{12}
\loigiai{
$6 \times 2 = 12$ cm.
}
\end{ex}

\begin{ex}
Một hình vuông có cạnh 6 cm được chia thành các ô vuông nhỏ cạnh 1 cm. Có tất cả bao nhiêu ô vuông nhỏ?
\shortans{36}
\loigiai{
$6 \times 6 = 36$ ô vuông.
}
\end{ex}

\begin{ex}
Một hình chữ nhật có chiều dài 5 cm, chiều rộng 3 cm được chia thành các ô vuông cạnh 1 cm. Có tất cả bao nhiêu ô vuông nhỏ?
\shortans{15}
\loigiai{
$5 \times 3 = 15$ ô vuông.
}
\end{ex}

\begin{ex}
Để gấp một chiếc thuyền giấy, bạn An dùng một tờ giấy hình vuông cạnh 20 cm. Chu vi tờ giấy đó là bao nhiêu xăng-ti-mét?
\shortans{80}
\loigiai{
$20 \times 4 = 80$ cm.
}
\end{ex}

\begin{ex}
Một họa tiết hình tam giác đều có cạnh dài 8 cm. Chu vi họa tiết đó là bao nhiêu xăng-ti-mét?
\shortans{24}
\loigiai{
$8 \times 3 = 24$ cm.
}
\end{ex}

\begin{ex}
Dùng 12 que tính bằng nhau xếp thành một hình vuông. Mỗi cạnh của hình vuông có bao nhiêu que tính?
\shortans{3}
\loigiai{
$12 : 4 = 3$ que tính.
}
\end{ex}

\begin{ex}
Dùng 9 que tính bằng nhau xếp thành 3 hình tam giác rời nhau. Mỗi tam giác có mấy que tính?
\shortans{3}
\loigiai{
$9 : 3 = 3$ que tính.
}
\end{ex}

\begin{ex}
Bạn Hoa cắt một hình chữ nhật có chu vi 30 cm thành hai hình vuông nhỏ bằng nhau cạnh 5 cm. Đúng hay sai? (Đúng ghi 1, sai ghi 0).
\shortans{1}
\loigiai{
Chu vi hình chữ nhật ghép từ 2 hình vuông cạnh 5 cm: chiều dài 10 cm, rộng 5 cm $\Rightarrow (10+5) \times 2 = 30$ cm.
}
\end{ex}

\begin{ex}
Bác thợ mộc đóng khung ảnh gỗ hình vuông cạnh 25 cm. Tổng chiều dài 4 thanh gỗ là bao nhiêu xăng-ti-mét?
\shortans{100}
\loigiai{
$25 \times 4 = 100$ cm.
}
\end{ex}

\begin{ex}
Để vẽ đường tròn có bán kính 8 cm, đường kính của nó là bao nhiêu xăng-ti-mét?
\shortans{16}
\loigiai{
$8 \times 2 = 16$ cm.
}
\end{ex}

\begin{ex}
Một bức tranh hình chữ nhật có chiều dài 40 cm, chiều rộng 30 cm. Nửa chu vi khung tranh là bao nhiêu xăng-ti-mét?
\shortans{70}
\loigiai{
$40 + 30 = 70$ cm.
}
\end{ex}

\begin{ex}
Một góc vuông được vẽ bằng ê-ke có số đo bao nhiêu độ? (Điền số).
\shortans{90}
\loigiai{
Góc vuông có số đo $90^\circ$.
}
\end{ex}

\begin{ex}
Có bao nhiêu góc vuông trong một hình vuông?
\shortans{4}
\loigiai{
Hình vuông có 4 góc vuông.
}
\end{ex}
